\documentclass[aps,prd,preprint,superscriptaddress,nofootinbib,longbibliography]{revtex4-2}

\usepackage[T1]{fontenc}
\usepackage{lmodern}
\usepackage{amsmath,amssymb,bm,mathtools}
\usepackage{booktabs,array,longtable}
\usepackage{xcolor}
\usepackage{hyperref}
\usepackage{microtype}

\hypersetup{
  colorlinks=true,
  linkcolor=blue!60!black,
  citecolor=blue!60!black,
  urlcolor=blue!60!black
}

% --- Macros ---
\newcommand{\dd}{\mathrm d}
\newcommand{\kt}{k_T}
\newcommand{\qT}{q_T}
\newcommand{\bT}{b_T}
\newcommand{\bb}{\bm b}
\newcommand{\bk}{\bm k}
\newcommand{\bq}{\bm q}
\newcommand{\bp}{\bm p}
\newcommand{\as}{\alpha_s}
\newcommand{\asbar}{a_s}
\newcommand{\KCSS}{\mathcal K_{\rm CSS2}}
\newcommand{\CSS}{\mathrm{CSS2}}
\newcommand{\ot}{\otimes_T}
\newcommand{\ox}{\otimes_x}
\newcommand{\deltaT}{\delta^{(2)}}
\newcommand{\NP}{\mathrm{NP}}
\newcommand{\pert}{\mathrm{pert}}
\newcommand{\FO}{\mathrm{FO}}
\newcommand{\ExpT}{\mathrm{Exp}_{\otimes_T}}
\newcommand{\NthreeLLp}{N^3\mathrm{LL}'}
\newcommand{\order}[1]{\mathcal O\!\left(#1\right)}
\newcommand{\mrm}[1]{\mathrm{#1}}
\newcommand{\tilF}{\widetilde F}
\newcommand{\tilC}{\widetilde C}
\newcommand{\tilD}{\widetilde D}

\begin{document}

\title{A pure transverse-momentum-space prescription for high-accuracy TMD extractions}

\author{Author list to be determined}
\affiliation{$k_T$ Perturbative Extraction Project}

\date{Draft version 0.1 -- June 15, 2026}

\begin{abstract}
We formulate a transverse-momentum-space representation of Collins/CSS2 transverse-momentum-dependent (TMD) factorization suitable for high-accuracy extractions without numerical transformations to impact-parameter space.  The construction, denoted $\KCSS$, uses two-dimensional plus distributions and transverse convolution kernels as the working objects.  The local Sudakov factor in $b_T$ space is replaced by a convolution-exponential Green function in $k_T$ space, while the TMD matching coefficients are imported by a coefficient-level map $L_b^m\mapsto \mathfrak L_{m-1}(\bk;\mu)$.  We define the perturbative accuracy label $N^k\mathrm{LL}'$ by ingredient equivalence to the parent CSS2 calculation: the same cusp and noncusp anomalous dimensions, Collins-Soper kernel, hard factors, TMD matching coefficients, DGLAP kernels, and fixed-order subtractions must be included, but represented in $k_T$ space.  The nonperturbative input is left as a model-agnostic module; neural networks are one possible parametrization, not part of the perturbative accuracy definition.  We also give the generic process-level $W$ operator, the $W+Y$ matching prescription, regulated normalization and moment constraints, and a validation hierarchy for future numerical implementations.  The result is a formal prescription for pure-$k_T$ TMD extractions capable of carrying a $\NthreeLLp$ perturbative label once the required physical three-loop coefficient payloads are supplied and validated.
\end{abstract}

\maketitle

\section{Introduction}

TMD factorization organizes the small-transverse-momentum region of color-singlet production and semi-inclusive reactions in terms of hard factors, TMD parton distributions or fragmentation functions, Collins-Soper evolution, and matching onto collinear parton distributions at short transverse distances~\cite{CollinsSoper1981,CSS1985,Collins2011,TMDHandbook2023}.  The standard high-accuracy implementation is formulated in impact-parameter space, where the transverse convolution becomes a product and the Sudakov factor is local in $b_T$.  This feature is technically powerful, and modern phenomenology has reached high logarithmic accuracy in $b_T$-space formulations~\cite{MoosScimemiVladimirovZurita2024,MoosScimemiVladimirovZurita2025,BillisMichelTackmann2025}.

There are nevertheless good reasons to develop an equivalent formulation that remains directly in transverse momentum.  First, momentum-space representations interface naturally with measured $q_T$ bins and with flexible nonperturbative models in $k_T$.  Second, a pure-$k_T$ implementation avoids fitting an object in one representation and repeatedly transforming it to another.  Third, recent neural-network and inverse-problem methods have shown that direct momentum-space reconstruction can be practical, although not yet with the perturbative content required for an $N^k\mathrm{LL}$ label~\cite{FernandoKeller2026,BacchettaEtAl2025NN}.  Finally, direct momentum-space and distribution-space resummation methods already demonstrate that logarithmic resummation need not be tied to numerical $b_T$ integration~\cite{EllisVeseli1998,MonniReTorrielli2016,BizonMonniReRottoliTorrielli2018,EbertTackmann2017,Simonelli2026}.

The purpose of this paper is to define a rigorous pure-$k_T$ prescription, not to report a fitted extraction.  We construct a $k_T$-space representation of the Collins/CSS2 scheme, denoted $\KCSS$, and show how the standard perturbative ingredients are represented by distribution-valued kernels and transverse convolution exponentials.  The central claim is deliberately limited:
\begin{quote}
A calculation in $\KCSS$ has the same $N^k\mathrm{LL}'$ perturbative accuracy as the corresponding CSS2 calculation when, and only when, the same perturbative ingredients are included and the fixed-order expansion is validated in $k_T$ space.
\end{quote}
Thus, the paper supplies the formal prescription needed for a future $\NthreeLLp$ extraction.  A numerical extraction at that accuracy additionally requires staged physical three-loop matching coefficients, code-level validation, data treatment, and a nonperturbative fit.

\section{Definition of the $\KCSS$ representation}

We define $\KCSS$ as the direct transverse-momentum-space representation of the CSS2 TMD scheme.  It is not introduced as a new operator scheme with new finite subtractions.  Instead,
\begin{equation}
  F_{i/h}^{\KCSS}(x,\bk;\mu,\zeta)
  \equiv
  \int \frac{\dd^2\bb}{(2\pi)^2}
  e^{-i\bb\cdot\bk}
  \tilF_{i/h}^{\CSS}(x,\bb;\mu,\zeta),
  \label{eq:Fmap}
\end{equation}
with analogous definitions for the Collins-Soper kernel and matching coefficients,
\begin{align}
  \mathcal D_i^{\KCSS}(\bk;\mu)
  &\equiv \mathcal F^{-1}_{b\to k}\!\left[\tilD_i^{\CSS}(b_T;\mu)\right],
  \\
  \mathcal C_{i\leftarrow j}^{\KCSS}(z,\bk;\mu,\zeta)
  &\equiv \mathcal F^{-1}_{b\to k}\!\left[\tilC_{i\leftarrow j}^{\CSS}(z,b_T;\mu,\zeta)\right].
\end{align}
Equation~\eqref{eq:Fmap} is a definition of the analytic distribution dictionary.  It is not a prescription to perform numerical Fourier-Bessel transforms during the fit.  Once the coefficient functions and evolution kernels are projected into the distribution basis, all subsequent operations are transverse convolutions in $k_T$.

The CSS2 $W$ term for a Drell-Yan-like process,
\begin{equation}
  W^{\CSS}(\bq)
  =\sum_{ij} H_{ij}(Q,\mu)
   \int \frac{\dd^2\bb}{(2\pi)^2}e^{-i\bb\cdot\bq}
   \tilF_{i/A}(x_A,\bb;\mu,\zeta_A)
   \tilF_{j/B}(x_B,\bb;\mu,\zeta_B),
\end{equation}
becomes
\begin{equation}
  W^{\KCSS}(\bq)
  =\sum_{ij} H_{ij}(Q,\mu)
  \left[F_{i/A}^{\KCSS}\ot F_{j/B}^{\KCSS}\right](\bq),
  \label{eq:Wkcss}
\end{equation}
where
\begin{equation}
  [A\ot B](\bq)=\int\dd^2\bk\,A(\bk)B(\bq-\bk).
\end{equation}
The hard factor is the parent CSS2 hard factor.  If a genuinely new $k_T$-adapted child scheme is later introduced, it must be related to $\KCSS$ by explicit finite transverse convolution kernels, not by an implicit reshuffling of terms.

\section{Distribution basis and convolution algebra}

The natural $k_T$-space objects are distributions.  We use the logarithmic two-dimensional plus distributions $\mathcal L_n(\bk;\mu)$, together with a derivative basis $\mathfrak L_n$ defined by its Fourier image,
\begin{equation}
  \mathcal F_{k\to b}\!\left[\mathfrak L_n(\bk;\mu)\right]
  =L_b^{n+1},
  \qquad
  L_b\equiv \ln\frac{\mu^2 b_T^2}{b_0^2},\quad b_0=2e^{-\gamma_E}.
  \label{eq:frak-def}
\end{equation}
The first relation to the conventional basis is
\begin{equation}
  \mathfrak L_0=-\mathcal L_0,
\end{equation}
with higher-order relations containing the corresponding zeta-valued contact terms.  The derivative basis has the simple convolution rule
\begin{equation}
  \boxed{\mathfrak L_m\ot \mathfrak L_n=\mathfrak L_{m+n+1}.}
  \label{eq:frak-algebra}
\end{equation}
Equivalently, if $\mathcal G_\eta$ is the analytic generator satisfying
\begin{equation}
  \mathcal F_{k\to b}\!\left[\mathcal G_\eta\right]=e^{\eta L_b},
\end{equation}
then
\begin{equation}
  \mathcal G_\eta\ot \mathcal G_\xi=\mathcal G_{\eta+\xi},
  \qquad
  \ExpT\{\eta\mathfrak L_0\}=\mathcal G_\eta .
\end{equation}
More generally,
\begin{equation}
  \mathcal F_{k\to b}\!\left[
  \ExpT\left\{\sum_{n\ge0}r_n\mathfrak L_n\right\}\right]
  =\exp\left[\sum_{n\ge0}r_n L_b^{n+1}\right].
  \label{eq:convolution-exponential-theorem}
\end{equation}
This identity is the core mechanism by which the local $b_T$-space Sudakov exponent is represented in $k_T$ space without numerical $b_T$ transformations.

For comparisons with data, distributions should be evaluated on test functions or bins, not pointwise at the origin.  For an annular bin $q_a<|\bq|<q_b$, define
\begin{equation}
  \mathcal B_n(q_a,q_b;Q)
  =\int_{q_a<|\bq|<q_b}\dd^2\bq\,\mathcal L_n(\bq;Q).
\end{equation}
For $q_a>0$ this reduces to the ordinary integral over the real-emission function, while the bin touching the origin contains the plus-prescription subtraction.  This bin-level treatment is essential for numerical stability and for comparison with experimental data.

\section{Evolution in $k_T$ space}

The CSS2 Collins-Soper equation becomes nonlocal in $k_T$ space,
\begin{equation}
  \frac{\partial}{\partial\ln\sqrt\zeta}
  F_{i/h}(x,\bk;\mu,\zeta)
  =-\mathcal D_i(\bk;\mu)\ot F_{i/h}(x,\bk;\mu,\zeta),
  \label{eq:cs-k-space}
\end{equation}
while the $\mu$ evolution remains multiplicative in transverse momentum,
\begin{equation}
  \frac{\partial}{\partial\ln\mu}
  F_{i/h}(x,\bk;\mu,\zeta)
  =\gamma_F^i(\mu,\zeta)F_{i/h}(x,\bk;\mu,\zeta).
\end{equation}
At a final scale $(\mu_f,\zeta_f)$ and initial scale $(\mu_i,\zeta_i)$, a convenient endpoint representation of the running-coupling Green function is
\begin{equation}
  \mathcal U_i(\bk;f,i)
  =\exp\!\left[R_i(\mu_f,\mu_i;\zeta_i)-\rho D_i^\delta(a_f)\right]
  \ExpT\!\left\{-\rho\sum_{n\ge0}D_{i,n}(a_f)\mathfrak L_n(\bk;\mu_f)\right\},
  \label{eq:running-U}
\end{equation}
where
\begin{equation}
  \rho=\ln\sqrt{\frac{\zeta_f}{\zeta_i}},\qquad a_f=\frac{\alpha_s(\mu_f)}{4\pi}.
\end{equation}
The scalar exponent is written in terms of the standard evolution integrals
\begin{equation}
  R_i=A_{\gamma_i}-\ell_i A_{\Gamma_i}+2S_{\Gamma_i},
  \qquad
  \ell_i=\ln\frac{\zeta_i}{\mu_i^2}.
\end{equation}
The coefficient functions $D_{i,n}(a_f)$ are determined by the perturbative Collins-Soper kernel in the derivative basis,
\begin{equation}
  \mathcal D_i(\bk;\mu)
  =D_i^\delta(a_\mu)\deltaT(\bk)+\sum_{n\ge0}D_{i,n}(a_\mu)\mathfrak L_n(\bk;\mu).
\end{equation}
Through two loops,
\begin{equation}
  \mathcal D_i(\bk;\mu)
  =a_s\Gamma_0^i\mathfrak L_0
  +a_s^2\left[\frac{1}{2}\beta_0\Gamma_0^i\mathfrak L_1+
  \Gamma_1^i\mathfrak L_0+d_i^{(2)}\deltaT\right]+
  \order{a_s^3}.
\end{equation}
The three-loop structure is fixed by the three-loop rapidity anomalous dimension together with the RG equation for $\mathcal D_i$; its explicit finite contact term is an imported perturbative ingredient at $\NthreeLLp$.

\section{Matching and the accuracy label}

At perturbative transverse momentum, the TMD matches onto collinear PDFs,
\begin{equation}
  F_{i/h}^{\pert}(x,\bk;\mu,\zeta)
  =\sum_j \mathcal C_{i\leftarrow j}^{\KCSS}(\bk;\mu,\zeta)\ox f_{j/h}(\mu).
  \label{eq:TMDmatching}
\end{equation}
The coefficient-level import rule is mechanical.  If the parent CSS2 coefficient is supplied as
\begin{equation}
  \tilC_{i\leftarrow j}^{(n),\CSS}(z,L_b,\ell_\zeta)
  =\sum_{m=0}^{2n}\tilde c_{i\leftarrow j}^{(n,m)}(z;\ell_\zeta)L_b^m,
  \qquad
  \ell_\zeta=\ln\frac{\zeta}{\mu^2},
\end{equation}
then
\begin{equation}
  \boxed{
  \mathcal C_{i\leftarrow j}^{(n),\KCSS}
  =\tilde c_{i\leftarrow j}^{(n,0)}\deltaT
  +\sum_{m=1}^{2n}\tilde c_{i\leftarrow j}^{(n,m)}\mathfrak L_{m-1}.}
  \label{eq:C-import}
\end{equation}
No numerical Fourier-Bessel transform appears in Eq.~\eqref{eq:C-import}.  At N$^3$LO the highest logarithmic power is $L_b^6$, so the $k_T$ basis must support $\mathfrak L_5$.

We define the $N^k\mathrm{LL}'$ label in $\KCSS$ by ingredient equivalence to the parent CSS2 calculation.  In the convention used here,
\begin{equation}
  \NthreeLLp_{\KCSS}
  =N^3\mathrm{LL}\ \text{evolution}
  +\text{three-loop hard and TMD matching boundary functions}.
\end{equation}
The required ingredients are summarized in Table~\ref{tab:ingredients}.  The prime is not a cosmetic label: it requires the fixed-order boundary functions at the same order as the claimed primed accuracy.

\begin{table}[t]
\caption{Perturbative ingredient manifest for a $\NthreeLLp_{\KCSS}$ singular $W$ operator in the convention used in this paper.  Equivalent conventions are possible, but the $k_T$ implementation must match the chosen parent CSS2 convention ingredient by ingredient.}
\label{tab:ingredients}
\begin{ruledtabular}
\begin{tabular}{ll}
Ingredient & Required order for $\NthreeLLp$ \\
\hline
Cusp anomalous dimension $\Gamma_{\rm cusp}$ & four loops \\
QCD beta function & four loops \\
TMD noncusp anomalous dimension $\gamma_F$ & three loops \\
Collins-Soper kernel / rapidity anomalous dimension & three loops \\
Hard factor $H$ & $\order{a_s^3}$ \\
TMD matching coefficients $\mathcal C_{i\leftarrow j}$ & $\order{a_s^3}$ \\
DGLAP kernels in boundary evolution & $P^{(0)},P^{(1)},P^{(2)}$ \\
$W$-term fixed-order expansion gate & through $\order{a_s^3}$ \\
$Y$ term for full-spectrum prediction & matching-consistent fixed order \\
\end{tabular}
\end{ruledtabular}
\end{table}

High-order ingredients entering this manifest are available in the literature: three-loop TMD matching coefficients for unpolarized TMDPDFs and TMDFFs~\cite{EbertMistlbergerVita2020,EbertMistlbergerVita2021FF,LuoYangZhuZhu2021}, rapidity anomalous dimensions through three loops and beyond~\cite{LiZhu2017,MoultZhuZhu2022,DuhrMistlbergerVita2022}, three-loop splitting functions~\cite{MochVermaserenVogt2004,VogtMochVermaseren2004}, and high-order hard functions from quark and gluon form factors~\cite{GehrmannGloverHuberIkizlerliStuderus2010}.  This paper does not print the long three-loop coefficient tables; they should be imported from source-stamped ancillary files and converted with Eq.~\eqref{eq:C-import}.

We distinguish a formal prescription from a numerical implementation.  A formal $\NthreeLLp_{\KCSS}$ prescription is specified by Eqs.~\eqref{eq:convolution-exponential-theorem}, \eqref{eq:running-U}, and \eqref{eq:C-import} plus the ingredient manifest.  A numerical implementation may claim $\NthreeLLp$ only after the physical three-loop payload is staged, converted, hashed, and validated by expansion through $\order{a_s^3}$.

\section{Process-level $W$ operator and full-spectrum matching}

The process-level singular operator is
\begin{equation}
  W_{AB\to F}^{\KCSS}(\bq;Q)
  =\sum_{ij}H_{ij\to F}(Q,\mu)
  \left[\mathbb F_{i/A}^{\KCSS}\ot \mathbb F_{j/B}^{\KCSS}\right](\bq),
  \label{eq:generic-W}
\end{equation}
with the appropriate replacement of one or both TMDPDFs by TMD fragmentation functions for SIDIS or $e^+e^-$ observables.  Examples include
\begin{align}
  \text{Drell-Yan:}&\quad F_{q/A}\ot F_{\bar q/B}+(q\leftrightarrow \bar q),\\
  \text{SIDIS:}&\quad F_{i/A}\ot D_{h/j},\\
  e^+e^-\text{ back-to-back hadrons:}&\quad D_{h_1/i}\ot D_{h_2/j},\\
  \text{gluon color singlets:}&\quad F_{g/A}\ot F_{g/B}.
\end{align}
Drell-Yan fixed-target data are therefore a useful first testbed, not a defining limitation of the formalism.

For a full transverse-momentum spectrum, the singular $W$ term must be matched to fixed order by a nonsingular remainder,
\begin{equation}
  \dd\sigma=\dd\sigma_W^{\KCSS}+\dd\sigma_Y,
  \qquad
  \dd\sigma_Y=\dd\sigma_{\FO}-\left[\dd\sigma_W^{\KCSS}\right]_{\FO}.
  \label{eq:WY}
\end{equation}
The TMD large-$k_T$ tail in Eq.~\eqref{eq:TMDmatching} is universal matching onto collinear PDFs or FFs.  The $Y$ term in Eq.~\eqref{eq:WY} is a cross-section-level fixed-order remainder.  These are distinct ingredients and should not be substituted for one another.  Fixed-order and matched codes such as DYTurbo, MCFM/CuTe-MCFM, SCETlib, or equivalent implementations can be used to benchmark or supply the fixed-order cross section~\cite{CamardaEtAlDYTurbo2020,CamardaCieriFerrera2021,NeumannCampbell2023,BillisMichelTackmann2025}.

\section{Model-agnostic nonperturbative module}

The perturbative prescription does not require a specific nonperturbative parametrization.  At a boundary scale $(\mu_T,\zeta_T)$, one may write schematically
\begin{equation}
  \mathbb B_{i/h}(x,\bk;\mu_T,\zeta_T)
  =\left[\mathcal C_{i\leftarrow j}\ox f_{j/h}\right](x,\bk)
  +\mathcal M_{i/h}^{\NP}(x,\bk;\theta)-\mathcal M_{i/h}^{\rm overlap}(x,\bk;\theta),
  \label{eq:NPadditive}
\end{equation}
or use a convolutional smearing form,
\begin{equation}
  \mathbb B_{i/h}
  =\left[\mathcal C_{i\leftarrow j}\ox f_{j/h}\right]
  \ot S_{i/h}^{\NP}+\text{power-suppressed remainder}.
  \label{eq:NPconv}
\end{equation}
Other choices, including nonperturbative modifications of the Collins-Soper kernel, are possible if they are specified in the chosen scheme and do not alter the perturbative ingredient manifest.

The nonperturbative module may be a neural network, spline basis, Gaussian mixture, analytic ansatz, lattice-informed input, or another parametrization.  The perturbative accuracy label belongs to the operator defined in the previous sections.  It does not depend on whether the nonperturbative module is a DNN.  Positivity, when imposed, should apply to physical cross sections or to well-defined low-scale model components, not to individual renormalized plus-distribution terms.

\section{Regulated normalization and moments}

A naive all-$k_T$ normalization condition,
\begin{equation}
  \int\dd^2\bk\,F_{i/h}(x,\bk;\mu,\zeta)=f_{i/h}(x,\mu),
\end{equation}
should not be imposed as an exact, unregulated identity once perturbative tails are present.  The large-$k_T$ tail produces ultraviolet-sensitive integrals and moments.  Instead, one should define regulated cumulants or smeared moments,
\begin{equation}
  \mathcal I_K[F_{i/h}]
  =\int\dd^2\bk\,w_K(k_T)F_{i/h}(x,\bk;\mu,\zeta),
\end{equation}
with matching relation
\begin{equation}
  \mathcal I_K[F_{i/h}]
  =\sum_j Z_{i\leftarrow j}(K;\mu,\zeta)\ox f_{j/h}(x,\mu)
  +\order{\left(\frac{\Lambda_{\rm QCD}}{K}\right)^p}.
\end{equation}
Similarly, transverse-momentum widths should be defined with cut or smeared weights,
\begin{equation}
  \langle k_T^2\rangle_K
  =\frac{\int\dd^2\bk\,w_K(k_T)k_T^2F(x,\bk)}
  {\int\dd^2\bk\,w_K(k_T)F(x,\bk)}.
\end{equation}
This prescription applies to both fixed-target and collider kinematics, with the scale $K$ and weights chosen to maintain perturbative control.

\section{Validation hierarchy and claim discipline}

A pure-$k_T$ implementation should pass a validation ladder before it is used for extraction:
\begin{enumerate}
  \item distribution algebra and bin-integral tests for $\mathcal L_n$ and $\mathfrak L_n$;
  \item one-loop singular-spectrum checks in the diagonal Drell-Yan channel;
  \item flavor-vector toy-PDF tests separating $x$-space and transverse plus prescriptions;
  \item evolution path-independence and semigroup tests for Eq.~\eqref{eq:running-U};
  \item coefficient-import tests for Eq.~\eqref{eq:C-import};
  \item fixed-order expansion checks through the order required by the claimed accuracy;
  \item comparison of binned singular spectra to the parent CSS2 implementation;
  \item full $W+Y$ matching checks against independent fixed-order or resummed benchmarks.
\end{enumerate}
A compact claim table is given in Table~\ref{tab:claims}.

\begin{table}[t]
\caption{Claim discipline for the formalism and future implementations.}
\label{tab:claims}
\begin{ruledtabular}
\begin{tabular}{ll}
Claim & Requirement \\
\hline
Formal $\NthreeLLp_{\KCSS}$ prescription & algebra, evolution, import map, manifest \\
Numerical $\NthreeLLp_{\KCSS}$ implementation & physical $a_s^3$ payload and expansion gate \\
$\NthreeLLp$ extraction & validated implementation plus data fit \\
DNN-based $\NthreeLLp$ extraction & same perturbative operator; DNN in $\mathcal M_{\NP}$ \\
\end{tabular}
\end{ruledtabular}
\end{table}

A companion validation suite was used in constructing this draft and is intended to be released as a public GitHub/Zenodo repository before submission.  In the present draft, the repository citation is a placeholder and the local source bundle contains the algebraic validators.  The public repository should record coefficient-file hashes, source metadata, and the expansion-gate status.

\section{Conclusions}

We have given a pure transverse-momentum-space prescription for CSS2-equivalent TMD factorization.  The formulation replaces the local $b_T$-space Sudakov factor by a transverse convolution exponential acting on $k_T$-space distributions.  Matching coefficients are imported by an analytic coefficient-level projection $L_b^m\mapsto\mathfrak L_{m-1}$, and the perturbative accuracy label is inherited by ingredient equivalence to the parent CSS2 calculation.

This framework separates the perturbative operator from the nonperturbative module.  Consequently, a DNN can be used for the intrinsic transverse-momentum input without being part of the definition of $N^3\mathrm{LL}'$ accuracy.  The same formalism also supports analytic models, splines, Gaussian mixtures, or lattice-informed inputs.  Future phenomenological work must stage the physical three-loop matching coefficients, pass the $\order{a_s^3}$ expansion gate, implement full $W+Y$ matching, and then perform the data extraction.  The formalism developed here is designed to be the perturbative backbone for that program.

\appendix

\section{One-loop diagonal check}

For bins away from $q_T=0$, the one-loop diagonal Drell-Yan singular structure obtained from the $k_T$ kernels is
\begin{equation}
  \pi a_s\left[4C_F\mathcal L_1(\bq;Q)-6C_F\mathcal L_0(\bq;Q)\right]
  =\frac{\alpha_s}{\pi}\frac{1}{q_T^2}
  \left[C_F\ln\frac{Q^2}{q_T^2}-\frac32C_F\right],
\end{equation}
which reproduces the standard $A_q^{(1)}=C_F$ and $B_q^{(1)}=-3C_F/2$ singular coefficients.  This check fixes the signs and normalizations of the $\mathcal L_0,\mathcal L_1$ basis and the conversion between $d^2\bq$ and $dq_T^2$ spectra.

\section{Companion-code placeholder}

The public repository entry below is intentionally a placeholder in this draft.  It should be replaced by a citable archived release before journal submission.

\begin{quote}
$k_T$ Perturbative Extraction Project, \emph{KCSS validation suite}, GitHub repository and Zenodo archive, URL and DOI to be supplied before submission.
\end{quote}

\begin{thebibliography}{99}

\bibitem{CollinsSoper1981}
J.~C. Collins and D.~E. Soper,
\textit{Back-to-back jets in QCD},
Nucl. Phys. B \textbf{193}, 381 (1981);
Erratum: Nucl. Phys. B \textbf{213}, 545 (1983).

\bibitem{CSS1985}
J.~C. Collins, D.~E. Soper, and G.~F. Sterman,
\textit{Transverse momentum distribution in Drell-Yan pair and W and Z boson production},
Nucl. Phys. B \textbf{250}, 199 (1985).

\bibitem{Collins2011}
J.~C. Collins,
\textit{Foundations of Perturbative QCD}
(Cambridge University Press, Cambridge, 2011).

\bibitem{TMDHandbook2023}
R.~Boussarie et al.,
\textit{TMD Handbook},
arXiv:2304.03302 [hep-ph].

\bibitem{MoosScimemiVladimirovZurita2024}
V.~Moos, I.~Scimemi, A.~Vladimirov, and P.~Zurita,
\textit{Extraction of unpolarized transverse momentum distributions from the fit of Drell-Yan data at N$^4$LL},
JHEP \textbf{05}, 036 (2024),
arXiv:2305.07473 [hep-ph].

\bibitem{MoosScimemiVladimirovZurita2025}
V.~Moos, I.~Scimemi, A.~Vladimirov, and P.~Zurita,
\textit{Determination of unpolarized TMD distributions from the fit of Drell-Yan and SIDIS data at N$^4$LL},
JHEP \textbf{11}, 134 (2025),
arXiv:2503.11201 [hep-ph].

\bibitem{BillisMichelTackmann2025}
G.~Billis, J.~K.~L. Michel, and F.~J. Tackmann,
\textit{Drell-Yan transverse-momentum spectra at N$^3$LL$'$ and approximate N$^4$LL with SCETlib},
JHEP \textbf{02}, 170 (2025),
arXiv:2411.16004 [hep-ph].

\bibitem{FernandoKeller2026}
I.~P. Fernando and D.~Keller,
\textit{Deep Neural Network extraction of Unpolarized Transverse Momentum Distributions},
Phys. Rev. D \textbf{113}, 096017 (2026),
arXiv:2510.17243 [hep-ph].

\bibitem{BacchettaEtAl2025NN}
A.~Bacchetta, V.~Bertone, C.~Bissolotti, M.~Cerutti, M.~Radici, S.~Rodini, and L.~Rossi,
\textit{Neural-network extraction of unpolarized transverse-momentum-dependent distributions},
Phys. Rev. Lett. \textbf{135}, 021904 (2025),
arXiv:2502.04166 [hep-ph].

\bibitem{EllisVeseli1998}
R.~K. Ellis and S.~Veseli,
\textit{W and Z transverse momentum distributions: resummation in $q_T$ space},
Nucl. Phys. B \textbf{511}, 649 (1998),
hep-ph/9706526.

\bibitem{MonniReTorrielli2016}
P.~F. Monni, E.~Re, and P.~Torrielli,
\textit{Higgs transverse-momentum resummation in direct space},
Phys. Rev. Lett. \textbf{116}, 242001 (2016),
arXiv:1604.02191 [hep-ph].

\bibitem{BizonMonniReRottoliTorrielli2018}
W.~Bizo\'n, P.~F. Monni, E.~Re, L.~Rottoli, and P.~Torrielli,
\textit{Momentum-space resummation for transverse observables and the Higgs $p_\perp$ at N$^3$LL+NNLO},
JHEP \textbf{02}, 108 (2018),
arXiv:1705.09127 [hep-ph].

\bibitem{EbertTackmann2017}
M.~A. Ebert and F.~J. Tackmann,
\textit{Resummation of transverse momentum distributions in distribution space},
JHEP \textbf{02}, 110 (2017),
arXiv:1611.08610 [hep-ph].

\bibitem{Simonelli2026}
A.~Simonelli,
\textit{Transverse momentum resummation and analytic continuation into the deep infrared},
Eur. Phys. J. C \textbf{86}, 291 (2026),
arXiv:2509.06843 [hep-ph].

\bibitem{EbertMistlbergerVita2020}
M.~A. Ebert, B.~Mistlberger, and G.~Vita,
\textit{Transverse momentum dependent PDFs at N$^3$LO},
JHEP \textbf{09}, 146 (2020),
arXiv:2006.05329 [hep-ph].

\bibitem{EbertMistlbergerVita2021FF}
M.~A. Ebert, B.~Mistlberger, and G.~Vita,
\textit{TMD fragmentation functions at N$^3$LO},
JHEP \textbf{07}, 121 (2021),
arXiv:2012.07853 [hep-ph].

\bibitem{LuoYangZhuZhu2021}
M.-x. Luo, T.-Z. Yang, H.~X. Zhu, and Y.~J. Zhu,
\textit{Unpolarized quark and gluon TMD PDFs and FFs at N$^3$LO},
JHEP \textbf{06}, 115 (2021),
arXiv:2012.03256 [hep-ph].

\bibitem{LiZhu2017}
Y.~Li and H.~X. Zhu,
\textit{Bootstrapping rapidity anomalous dimensions for transverse-momentum resummation},
Phys. Rev. Lett. \textbf{118}, 022004 (2017),
arXiv:1604.01404 [hep-ph].

\bibitem{MoultZhuZhu2022}
I.~Moult, H.~X. Zhu, and Y.~J. Zhu,
\textit{The four loop QCD rapidity anomalous dimension},
JHEP \textbf{08}, 280 (2022),
arXiv:2205.02249 [hep-ph].

\bibitem{DuhrMistlbergerVita2022}
C.~Duhr, B.~Mistlberger, and G.~Vita,
\textit{Four-loop rapidity anomalous dimension and event shapes to fourth logarithmic order},
Phys. Rev. Lett. \textbf{129}, 162001 (2022),
arXiv:2205.02242 [hep-ph].

\bibitem{MochVermaserenVogt2004}
S.~Moch, J.~A.~M. Vermaseren, and A.~Vogt,
\textit{The three-loop splitting functions in QCD: The non-singlet case},
Nucl. Phys. B \textbf{688}, 101 (2004),
hep-ph/0403192.

\bibitem{VogtMochVermaseren2004}
A.~Vogt, S.~Moch, and J.~A.~M. Vermaseren,
\textit{The three-loop splitting functions in QCD: The singlet case},
Nucl. Phys. B \textbf{691}, 129 (2004),
hep-ph/0404111.

\bibitem{GehrmannGloverHuberIkizlerliStuderus2010}
T.~Gehrmann, E.~W.~N. Glover, T.~Huber, N.~Ikizlerli, and C.~Studerus,
\textit{Calculation of the quark and gluon form factors to three loops in QCD},
JHEP \textbf{06}, 094 (2010),
arXiv:1004.3653 [hep-ph].

\bibitem{CamardaEtAlDYTurbo2020}
S.~Camarda et al.,
\textit{DYTurbo: Fast predictions for Drell-Yan processes},
Eur. Phys. J. C \textbf{80}, 251 (2020),
arXiv:1910.07049 [hep-ph].

\bibitem{CamardaCieriFerrera2021}
S.~Camarda, L.~Cieri, and G.~Ferrera,
\textit{Drell-Yan lepton-pair production: $q_T$ resummation at N$^3$LL accuracy and fiducial cross sections at N$^3$LO},
Phys. Rev. D \textbf{104}, L111503 (2021),
arXiv:2103.04974 [hep-ph].

\bibitem{NeumannCampbell2023}
T.~Neumann and J.~Campbell,
\textit{Fiducial Drell-Yan production at the LHC improved by transverse-momentum resummation at N$^4$LL+N$^3$LO},
Phys. Rev. D \textbf{107}, L011506 (2023),
arXiv:2207.07056 [hep-ph].

\end{thebibliography}

\end{document}
